\section{Measuring the Task Value of an Upgrade}
\label{sec:method}

\subsection{Setting}
\label{sec:setting}

\paragraph{Tasks and agents.} We use \tautwo{} in text mode: a development panel of 58 retail and 27 airline tasks plus 80 telecom development tasks, and a holdout of 34 telecom tasks that no agent ran before its protocol was frozen. % src: manifests/fork/splits_v1.json
The agents are Qwen3-4B-Instruct (q4) and Qwen3-8B with thinking disabled (q8); a second family, Mistral-Small-3.2-24B, is a scope check (Table~\ref{tab:scope}). % nocheck: model name
One greedy Qwen3-30B-A3B-Instruct user simulator serves all agents. % src: ../cloq_work/sims/*.json (hf_id); docs/FINAL_EXPERIMENTS.md
Retail and airline success is the final-database match with normal termination, without all of \tautwo{}'s official communication checks; telecom uses the official task-basis success. BFCL multi-turn and ALFWorld run once per item without a user simulator (Appendix~\ref{app:setup}).

\paragraph{Quantization and serving.} Q is a simulated (fake-quantized) checkpoint served with BF16 compute on the same path as B. The main configurations use GPTQ \citep{frantar2023gptq} and AWQ-style scaling \citep{lin2024awq} at W4 (group size 128), calibrated on BF16 agent trajectories of 57 retail and airline tasks disjoint from the development tasks, and an INT8 KV cache (KV8). Stress configurations (GPTQ-W3, RTN-W3, rotated KV4) and QuaRot, SpinQuant and KIVI \citep{ashkboos2024quarot,liu2025spinquant,liu2024kivi} are scope checks (Table~\ref{tab:scope}; Appendix Table~\ref{tab:qconfig}). % src: runs/a6k2_q4_gptq_w4_d1_f01/quantization.json; runs/a6k2_q8_awq_w4_f01/quantization.json; runs/a6k2_q8_kvq_int8_f01/quantization.json; manifests/a6k2/q4/calib_fp_traj_c1.json
B and Q run on vLLM \citep{kwon2023vllm} with batch-invariant kernels \citep{he2025defeating} and deterministic compilation in separate server processes (no shared prefix cache). A \emph{block} is one closed-loop run of Q over the task pool under one prompt nonce (a random session-identifier line in the system prompt); blocks are the replicates, and a \emph{unit} is one task in one block.

\subsection{Four branches on one state}
\label{sec:estimands}

\paragraph{Frozen state.} Before Q's decision $t$ we restore the state $s_t$: the full message history; the environment and user-device databases, rebuilt by strict replay of every recorded tool call; the user simulator's visible conversation; and the used step and error budget. B and Q act from identical prompt token ids, and the state is written to disk before any branch runs.

\paragraph{Four branches.} $\YY{XY}\in\{0,1\}$ is task success when the action at $s_t$ comes from X and model Y executes the rest of the episode, X, Y $\in\{\mathrm{B},\mathrm{Q}\}$; all four branches share the remaining budget, and each upgrade option is a paired contrast with QQ: $G$ (this step only), $H$ (from the next step) and $F$ (both), defined in Figure~\ref{fig:case}. % nocheck: notation {0,1}
A \emph{rescue} is a unit where the upgraded branch succeeds and QQ fails, a \emph{damage} the reverse; \emph{turnover} is rescues plus damages, \emph{net} rescues minus damages. An outcome is \emph{continuation dependent} when the action swap's effect differs between Q's and B's continuation. All contrasts measure effects on task success, not which action is correct (symbols: Appendix~\ref{app:names}).

\paragraph{Reference interventions.} A large turnover could mean that a decision needs higher precision, or only that it is sensitive to a change of model. We therefore rerun the design with Q$'$ in place of B: a second GPTQ-W4 checkpoint of the same model, calibrated on a re-drawn state and row sample of the same calibration trajectories, with the same Q episodes, tasks and nonce blocks and a protocol frozen before any outcome (q4 and q8 GPTQ-W4; GPTQ-W3 as a stress check). \S\ref{sec:sameprec} adds common states and a key-order change of Q's own action.

% Method schematic (R2 review): states and denominators, interval conventions, replay acceptance.
% The numbers are set as macros on their own lines so that each carries its "% src:" (TikZ layout lines are nocheck).
\newcommand{\gateQQ}{20/20}% src: docs/insights/FORK_PROTOCOL.md §8 (analysis/fork/gate/)
\newcommand{\gateBB}{10/10}% src: docs/insights/FORK_PROTOCOL.md §8 (analysis/fork/gate/)
\newcommand{\noopdev}{924/924}% src: analysis/fork/out/dev_all_20261001T070326Z.md (TL;DR: No-op sample 924/924)
\newcommand{\noopfrac}{10\%}% src: docs/insights/FORK_PROTOCOL.md §9 (QQ no-op on a 10% hash sample; stop rule)
\newcommand{\minblk}{10}% nocheck: interval rule as in archived method L39 (block bootstrap with at least 10 blocks)
\begin{figure}[t]
\centering
\begin{tikzpicture}[ % nocheck: TikZ layout
  font=\scriptsize, % nocheck: TikZ layout
  hd/.style={font=\scriptsize\bfseries,anchor=north west,inner sep=0pt}, % nocheck: TikZ layout
  txt/.style={anchor=north west,align=left,inner sep=0pt,text width=5.15cm}, % nocheck: TikZ layout
  st/.style={circle,inner sep=0pt,minimum size=4pt,fill=black!25}] % nocheck: TikZ layout
% (a) states and denominators
\node[hd] at (0,0) {(a) Which states, which denominator}; % nocheck: TikZ layout
\node[anchor=north west,inner sep=0pt,text=black!60] at (0,-.32) {one task$\times$block unit: Q's episode, decisions $t$}; % nocheck: TikZ layout
\draw[-{Stealth[length=3pt]},black!45] (.1,-1.0) -- (7.55,-1.0) node[right,inner sep=1pt,text=black!60] {$t$}; % nocheck: TikZ layout
\foreach \x in {.5,1.3,2.1,2.9,3.7,4.5,5.3,6.1,6.9} \node[st] at (\x,-1.0) {}; % nocheck: TikZ layout
\node[circle,inner sep=0pt,minimum size=6pt,fill=palDamage,draw=white] (fd) at (2.1,-1.0) {}; % nocheck: TikZ layout
\node[diamond,inner sep=0pt,minimum size=7pt,fill=palB,draw=white] (os) at (5.3,-1.0) {}; % nocheck: TikZ layout
\node[anchor=south,align=center,inner sep=1pt,text=palDamage] at (2.1,-.88) {\textbf{first divergent state}}; % nocheck: TikZ layout
\node[anchor=south,align=center,inner sep=1pt,text=palB] at (5.3,-.88) {\textbf{ordinary sampled state}}; % nocheck: TikZ layout
\node[anchor=north,align=center,inner sep=1pt,text width=3.5cm] at (2.1,-1.13) {first decision where B's and Q's actions differ in kind, tool, call count, order or normalised arguments (not wording)}; % nocheck: TikZ layout
\node[anchor=north,align=center,inner sep=1pt,text width=3.4cm] at (5.65,-1.13) {one per unit by a hash fixed before generation (Q's or BF16 trajectories); a \emph{common state} when several action sources act on it}; % nocheck: TikZ layout
\node[draw=black!25,fill=gray!6,rounded corners=1.5pt,anchor=north west,align=left,text width=3.45cm,inner sep=2.5pt] at (0,-2.62) {\textbf{full pool}: rescues, damages, turnover, net per task$\times$block unit; no divergence counts 0}; % nocheck: TikZ layout
\node[draw=black!25,fill=gray!6,rounded corners=1.5pt,anchor=north west,align=left,text width=3.45cm,inner sep=2.5pt] at (3.85,-2.62) {\textbf{divergent states}: $G$, $H$, $F$, continuation dependence, per configuration}; % nocheck: TikZ layout
% (b) interval conventions
\node[hd] at (0,-3.85) {(b) Interval conventions (Figs.~\ref{fig:rd}--\ref{fig:refs})}; % nocheck: TikZ layout
\draw[black!80,line width=.6pt] (.15,-4.55) -- (1.85,-4.55); \draw[black!80,line width=1.8pt] (.5,-4.55) -- (1.5,-4.55); \fill[black!80] (1.0,-4.55) circle (1.6pt); % nocheck: TikZ layout
\node[txt] at (1.95,-4.2) {thick: \textbf{fixed pool} 95\%, primary (block bootstrap, ${\geq}\minblk$ blocks, else $t$ over block means); thin: \textbf{new tasks} 95\% (crossed task$\times$block bootstrap)}; % nocheck: interval rule as in archived method L39
\draw[black!80,line width=1.0pt] (.3,-5.5) -- (1.7,-5.5); \fill[black!80] (.75,-5.5) circle (1.6pt); % nocheck: TikZ layout
\node[txt] at (1.95,-5.3) {telecom holdout: one crossed 90\% bar, the frozen reading interval}; % nocheck: interval rule (PREREG_FORK_CONFIRM)
\draw[black!80,line width=1.8pt] (.3,-6.16) -- (1.7,-6.16); \fill[black!80] (1.0,-6.16) circle (1.6pt); \draw[black!80,line width=.6pt] (.45,-6.32) -- (1.55,-6.32); \draw[black!80,line width=.6pt] (.45,-6.27) -- (.45,-6.37) (1.55,-6.27) -- (1.55,-6.37); % nocheck: TikZ layout
\node[txt] at (1.95,-5.98) {common states: thick two-look 97\% (the reading), thin unadjusted 95\% on its own sub-row}; % nocheck: interval rule (COMMONSTATE_SOURCES_PROTOCOL D2.2)
\draw[black!80,line width=1.0pt] (.3,-6.85) -- (1.7,-6.85); \fill[black!80] (1.0,-6.85) circle (1.6pt); % nocheck: TikZ layout
\node[txt] at (1.95,-6.68) {key order: one bar, wider of block / crossed 95\%}; % nocheck: interval rule (KEYORDER_PERTURB_PROTOCOL)
\node[anchor=north west,inner sep=0pt,text width=7.4cm,align=left,text=black!70] at (0,-7.15) {\emph{Resolved}: the named interval excludes zero; \emph{not resolved}: the reported interval(s) include zero. Fig.~\ref{fig:trigger} reports exact one-sided $p$-values, no intervals.}; % nocheck: TikZ layout
% (c) replay acceptance card
\node[draw=palRescue,fill=palRescue!8,rounded corners=2pt,anchor=north west,align=left,text width=7.25cm,inner sep=3pt] at (0,-8.2) {\textbf{(c) Replay acceptance.} Gate, before any development outcome: no-op replays Q$\rightarrow$Q \gateQQ{} and B$\rightarrow$B \gateBB{} states identical in prompt-token hashes, messages, database hashes, termination, reward and budget counters. Production: QQ no-op on a \noopfrac{} hash sample of every fork task's states (one mismatch quarantines the task): \noopdev{} sampled development states passed.}; % nocheck: TikZ layout
\end{tikzpicture}
\caption{Method at a glance. (a) The two state sources and the denominator each measure uses. (b) Interval conventions used in Figures~\ref{fig:rd} and~\ref{fig:refs}; Figure~\ref{fig:trigger} gives exact $p$-values. (c) Replay acceptance (Appendix~\ref{app:setup}).}
\label{fig:method}
\end{figure}

\subsection{Making the counterfactuals auditable}
\label{sec:replay}

\paragraph{Replay acceptance.} Before any development outcome was analysed, exact no-op replays on real divergence states passed a gate, and in production every fork task re-ran the QQ no-op on an outcome-independent hash sample of its states; all sampled states of every state source and of the same-precision run passed (Figure~\ref{fig:method}c; Appendix~\ref{app:setup}). % src: docs/insights/FORK_PROTOCOL.md §8, §9; analysis/fork/out/dev_all_20261001T070326Z.md
Exact replay is a tested property of this stack (greedy decoding, vLLM 0.29.0, batch-invariant kernels, deterministic compilation, dense BF16 kernels), not a general one \citep{yuan2025nondeterminism}; determinism tests and untested conditions are in Appendix~\ref{app:setup}. % src: analysis/status/gpu_fill_20261003/rerun_premise_det.json; docs/RESEARCH_LOG.md (W100w determinism-premise table); analysis/fork/commonstate_replay_check/REPLAY_8of9.md (TL;DR; serving context)

\subsection{Which states}
\label{sec:states}

\paragraph{State sources.} A \emph{common state} is an ordinary sampled state at which several action sources are compared; Figure~\ref{fig:method}a gives both state sources and the denominator of each measure. % src: docs/insights/FORK_PROTOCOL.md §15 amendment 1
Contrasts are paired within a configuration and descriptive across configurations, which select their own divergences; only complete blocks enter. % nocheck: interval rule as in archived method L39
